\begin{problem}[Minimal primes as component-generic points]
Prove that the irreducible components of $\Spec A$ are exactly
\[
V(\mathfrak p),\qquad \mathfrak p\in\operatorname{Min}(A),
\]
and that each minimal prime $\mathfrak p$ is the unique generic point of its component.
\end{problem}

\paragraph{Why this problem matters.}
It turns primary decomposition intuition into a geometric statement about distinguished
points.

\paragraph{Useful ideas before starting.}
Every irreducible closed subset is $V(\mathfrak p)$ for a prime.  Inclusion of $V$-sets
reverses inclusion of ideals.

\paragraph{Guided hints.}
Ask when $V(\mathfrak p)$ is maximal among irreducible closed subsets.

\begin{solution}
Every irreducible closed subset of $\Spec A$ is $V(\mathfrak p)$ for a prime
$\mathfrak p$.  It is an irreducible component exactly when it is maximal under inclusion
among irreducible closed subsets.

If $\mathfrak q\subsetneq\mathfrak p$, then
\[
V(\mathfrak p)\subsetneq V(\mathfrak q).
\]
Therefore $V(\mathfrak p)$ is maximal irreducible exactly when there is no smaller prime
$\mathfrak q\subsetneq\mathfrak p$, that is, exactly when $\mathfrak p$ is minimal.
Finally,
\[
\overline{\{\mathfrak p\}}=V(\mathfrak p),
\]
so $\mathfrak p$ is the component's unique generic point.
\end{solution}

\paragraph{What happened mathematically.}
Maximal geometric pieces correspond to minimal prime ideals because geometry reverses ideal
containment.

\paragraph{Extensions.}
Apply the theorem to $k[x,y]/(xy)$ and $k[x,y,z]/(xy,xz)$.

\paragraph{Provenance.}
New canonical problem; the minimal-prime/component relationship was already developed in
VI/05 and is now recast pointwise.
